\kbvsechead{(3)\hspace{0.45em}ಯಮಕಾಲಂಕಾರ :}

\indent ಒಂದು ಪದ್ಯದಲ್ಲಿ ಮೂರು ಅಥವಾ ಹೆಚ್ಚು ಅಕ್ಷರಗಳಿಂದ ಕೂಡಿದ ಪದವು ನಿಯತವಾಗಿ ಬಂದಿದ್ದರೆ, ಅದು ಯಮಕಾಲಂಕಾರವೆನಿಸುವುದು.\par

\noindent\hspace*{1.2em}\kbvlabel{ಉದಾ}\hspace{0.6em}ಯತಿ ಕೆಡುಗು ದುಸ್ಸಂಗದಲಿ ಭೂ ಪತಿ ಕೆಡುಗು ದುರ್ಮಂತ್ರಿಯಲಿ ನರ ಸುತ ಕೆಡುಗು ಲಾಲನೆಗಳಲಿ ಕೃಷಿ ಕೆಡುಗುಪೇಕ್ಷೆಯಲಿ ಮತಿ ಕೆಡುಗು ಮಧುಪಾನದಲಿ ಸ ದ್ಗತಿ ಕೆಡುಗು ದುಶ್ಚರಿತದಲಿ ನಿಜ\par

\begin{kbvverse}
\noindent ಸತಿ ಕೆಡುಗು ದುರ್ವ್ಯಸನದಲಿ ಭೂಪಾಲ ಕೇಳೆಂದ ॥
\end{kbvverse}

\noindent\hfill{\fontsize{9.5pt}{12.5pt}\selectfont\color{kbvgrey}—ಕರ್ಣಾಟ ಭಾರತ ಕಥಾ ಮಂಜರಿ : ಉದ್ಯೋಗ ಪರ್ವ}\par\vspace{3pt plus 2pt}

\indent ಇಲ್ಲಿ “ಕೆಡುಗು” ಎಂಬ ಶಬ್ದ ಪದ್ಯದ ಪ್ರತಿಸಾಲಿನಲ್ಲಿ ಎರಡನೆಯ ಪದವಾಗಿ ಪುನರುಕ್ತವಾಗಿದ್ದು, ಆ ಪದ್ಯದ ಸೊಬಗನ್ನು ಹೆಚ್ಚಿಸಿದೆ. “ಕೆಡುಗು” ಪದದಲ್ಲಿ ಮೂರು ಅಕ್ಷರಗಳಿರುವ ಕಾರಣ ಇದು ಯಮಕಾಲಂಕಾರ ವೆನಿಸುವುದು.\par

\kbvexercise{ಅಭ್ಯಾಸ ಪತ್ರಿಕೆ ೨೦}

\vspace{2.5pt plus 3pt minus 0.5pt}\noindent\hspace*{1.2em}\kbvnum{(1)}\hspace{0.55em}ಭಾಷಾಲಂಕಾರದ ಎರಡು ಶಾಖೆಗಳು ಯಾವುವು?\par

\vspace{2.5pt plus 3pt minus 0.5pt}\noindent\hspace*{1.2em}\kbvnum{(2)}\hspace{0.55em}ಅರ್ಥಾಲಂಕಾರ ಎಂದರೇನು? ಅತ್ಯಂತ ಸಮಂಜಸವಾದ ಉದಾಹರಣೆ ಸಹಿತ ವಿವರಿಸಿ.\par

\vspace{2.5pt plus 3pt minus 0.5pt}\noindent\hspace*{1.2em}\kbvnum{(3)}\hspace{0.55em}ಶಬ್ದಾಲಂಕಾರ ಎಂದರೇನು? ಅತ್ಯಂತ ಸಮಂಜಸವಾದ ಉದಾಹರಣೆ ಸಹಿತ ವಿವರಿಸಿ.\par

\vspace{2.5pt plus 3pt minus 0.5pt}\noindent\hspace*{1.2em}\kbvnum{(4)}\hspace{0.55em}ಅರ್ಥಾಲಂಕಾರದ ಯಾವುದಾದರೂ ಎರಡು ಪ್ರಕಾರಗಳನ್ನು ಹೆಸರಿಸಿ.\par

\vspace{2.5pt plus 3pt minus 0.5pt}\noindent\hspace*{1.2em}\kbvnum{(5)}\hspace{0.55em}ಶಬ್ದಾಲಂಕಾರದ ಯಾವುದಾದರೂ ಎರಡು ಪ್ರಕಾರಗಳನ್ನು ಹೆಸರಿಸಿ.\par
